\documentclass[12pt,a4paper]{article}
\usepackage[utf8]{inputenc}
\usepackage[T1,T2A]{fontenc}
\usepackage{amsmath,amssymb,amsfonts,mathtools}
\usepackage{amsthm}
\usepackage{geometry}
\geometry{left=1.0cm,right=1.0cm,top=1.25cm,bottom=3.1cm}
\usepackage{booktabs}
\usepackage{tabularx}
\usepackage{array}
\usepackage{hyperref}
\usepackage{url}
\usepackage{graphicx}
\usepackage{xcolor}
\usepackage{titlesec}
\usepackage{fancyhdr}
\usepackage{microtype}
\usepackage{multicol}
\usepackage{fontspec}
\usepackage{polyglossia}
\setdefaultlanguage{russian}
\setotherlanguage{english}
\usepackage{tikz}
\usepackage{pgfplots}
\pgfplotsset{compat=1.18}
\usetikzlibrary{positioning, arrows.meta, shapes.geometric, calc}
\usepackage{enumitem}

% ========= НАСТРОЙКА ШРИФТОВ ДЛЯ РУССКОГО ЯЗЫКА (Windows) =========
\defaultfontfeatures{Ligatures=TeX}
\setmainfont{Times New Roman}[
Script=Cyrillic,
Language=Russian
]
\setsansfont{Arial}[
Script=Cyrillic,
Language=Russian,
BoldFont=Arial Bold,
ItalicFont=Arial Italic,
BoldItalicFont=Arial Bold Italic
]
\setmonofont{Courier New}[
Script=Cyrillic,
Language=Russian,
BoldFont=Courier New Bold,
ItalicFont=Courier New Italic,
BoldItalicFont=Courier New Bold Italic
]
% Обязательное определение для polyglossia (моноширинный шрифт для кириллицы)
\newfontfamily\cyrillicfonttt{Courier New}[
Script=Cyrillic,
Language=Russian,
BoldFont=Courier New Bold,
ItalicFont=Courier New Italic,
BoldItalicFont=Courier New Bold Italic
]
% ==================================================================

% Цветовая схема YUCT
\definecolor{skyblue}{RGB}{0,102,204}
\definecolor{darkblue}{RGB}{0,0,139}
\definecolor{gold}{RGB}{255,215,0}

% Макросы YUCT
\newcommand{\Keff}{K_{\mathrm{eff}}}
\newcommand{\kappac}{\kappa_c}
\newcommand{\eps}{\varepsilon}
\newcommand{\betaf}{\beta}
\newcommand{\df}{d_{\!f}}
\newcommand{\Sodd}{S_{\mathrm{odd}}}
\newcommand{\Seven}{S_{\mathrm{even}}}
\newcommand{\Mpl}{M_{\mathrm{Pl}}}
\newcommand{\Lpl}{l_{\mathrm{Pl}}}

% Теорематические окружения (на русском)
\newtheorem{theorem}{Теорема}[section]
\newtheorem{axiom}[theorem]{Аксиома}
\newtheorem{remark}[theorem]{Замечание}
\newtheorem{definition}[theorem]{Определение}
\renewcommand{\proofname}{Доказательство}

% Форматирование заголовков разделов
\titleformat{\section}{\LARGE\bfseries\color{darkblue}}{\thesection}{1em}{}
\titleformat{\subsection}{\Large\bfseries\color{darkblue}}{\thesubsection}{1em}{}

% Колонтитулы для первой страницы
\fancypagestyle{firstpage}{%
	\fancyhf{}%
	\renewcommand{\headrulewidth}{0pt}%
	\renewcommand{\footrulewidth}{0pt}%
	\fancyfoot[C]{%
		\begin{minipage}[t]{\textwidth}
			\centering
			\textcolor{gold}{\rule{\textwidth}{1.6pt}}\\[5pt]
			{\small \textsuperscript{1}Yakushev Research, \url{https://yuct.org/}} \hfill
			{\small Протокол YPSDC: \url{https://ypsdc.com/}}\\[-2pt]
			\textcolor{gold}{\rule{\textwidth}{1.6pt}}\\[5pt]
			\textbf{ЕДИНАЯ ТЕОРИЯ КООРДИНАЦИИ ЯКУШЕВА 2026} \hfill \thepage\\[10pt]
		\end{minipage}%
	}%
}

\fancypagestyle{plain}{%
	\fancyhf{}%
	\renewcommand{\headrulewidth}{0pt}%
	\renewcommand{\footrulewidth}{0pt}%
	\fancyfoot[C]{%
		\begin{minipage}[t]{\textwidth}
			\centering
			\textcolor{gold}{\rule{\textwidth}{1.6pt}}\\[4pt]
			\textbf{ЕДИНАЯ ТЕОРИЯ КООРДИНАЦИИ ЯКУШЕВА 2026} \hfill \thepage\\[4pt]
		\end{minipage}%
	}%
}

\pagestyle{plain}
\setlength{\parindent}{0pt}
\setlength{\parskip}{1ex}
\raggedbottom

\hypersetup{
	colorlinks=true,
	linkcolor=blue,
	citecolor=blue,
	urlcolor=blue,
}

% Пользовательский макрос для шапки
\newcommand{\makeheader}{%
	\noindent
	\begin{minipage}{0.17\textwidth}
		\raggedright
		{\sffamily\bfseries\color{skyblue}\fontsize{46}{46}\selectfont YUCT}%
	\end{minipage}%
	\hspace{30pt}
	\begin{minipage}{0.32\textwidth}
		\raggedright
		{\sffamily\fontsize{11}{13}\selectfont ЕДИНАЯ \\ ТЕОРИЯ КООРДИНАЦИИ\\ ЯКУШЕВА}%
	\end{minipage}%
	\hfill
	\begin{minipage}{0.38\textwidth}
		\raggedleft
		{\sffamily\bfseries Первые принципы}\\ 
		{\sffamily Версия 7.0 / Май 2026}\\ 
		{\sffamily\footnotesize \scriptsize \url{https://doi.org/10.5281/zenodo.18444598}}%
	\end{minipage}
	\par\vspace{5pt}
	\noindent\textcolor{gold}{\rule{\textwidth}{1.6pt}}
	\par\vspace{0pt}
}

\begin{document}
	
	% Титульная страница
	\thispagestyle{firstpage}
	\makeheader
	
	\vspace{0cm}
	{\LARGE\bfseries Первые принципы и полный вывод YUCT \\
		\Large Вселенная как сеть координации: от аксиом к \\ живой клетке и \(\alpha \approx 1/137\) \par}
	\vspace{0.2cm}
	
	{\large Алексей В. Якушев\textsuperscript{1}} \\
	\vspace{0.0cm}
	
	{\color{darkblue}\textbf{\LARGE Аннотация}} \par
	{\large
		\noindent
		Мы представляем полное аксиоматическое основание Единой теории координации Якушева (YUCT). Реальность представляет собой иерархическую сеть координации, управляемую онтологическим первым принципом \(K_{\mathrm{eff}} > 1\). Из единственной аксиомы (фрактальная масштабная инвариантность) и определений протокола YPSDC мы доказываем три теоремы: минимальная энтропия словаря равна точно двум битам, пространственная размерность необходимо равна трём, а время является эмерджентной мерой несовершенства координации. Универсальный закон ошибки \(\varepsilon \propto K_{\mathrm{eff}}^{-2/3}\) устанавливается как эмпирически подтверждённый закон, и из него выводится алгебраический цикл \(S_{\mathrm{odd}} = 1,2\), \(S_{\mathrm{even}} = 0,8\). Масштабный квант \(q = (3/2)^{1/3}\) порождает универсальную массовую лестницу от планковской шкалы до живых клеток. Постоянная тонкой структуры \(\alpha^{-1} \approx 137,036\) получается из топологического инварианта \(N_{\mathrm{gauge}} = 12\) на компактном слое \(F_{18}\) с универсальной поправкой \(\beta\gamma = 0,20\). Космологическая постоянная \(\Omega_\Lambda = 2/3\) следует из той же топологии. Электрослабая шкала возникает из подсчёта фермионов Вейля \(L = 96\). Мы явно различаем теоремы, эмпирические законы и феноменологические подгонки: полуцелые индексы \(N_f\) для фермионов и фактор перекрытия \(W\)-бозона \(\xi_W \approx 0,53\) в настоящее время являются феноменологическими параметрами, чьё топологическое происхождение остаётся открытой проблемой. Теория не содержит подстраиваемых непрерывных параметров и даёт несколько конкретных, фальсифицируемых предсказаний в физике, биологии и космологии.
	}
	
	\vspace{0.1cm}
	\noindent
	{\color{darkblue}\textbf{Ключевые слова:}} YUCT, эффективность координации, алгебраический цикл, массовая лестница, постоянная тонкой структуры, эмерджентная геометрия, квантование времени, космологическая постоянная, фальсифицируемость.
	\vspace{0.1cm}
	\noindent
	{\color{darkblue}\textbf{Цитирование:}} Якушев А.В. Первые принципы и полный вывод YUCT. 2026. DOI: \url{https://doi.org/10.5281/zenodo.18444598} (настоящий том).
	
	\vspace{0.4cm}
	
	% ========================
	% РАЗДЕЛ 1: ОНТОЛОГИЯ И ОПРЕДЕЛЕНИЯ
	% ========================
	\section{Онтология и определения: координация как первичное понятие}
	\label{sec:ontology}
	
	Единая теория координации Якушева (YUCT) основана на единственном онтологическом примитиве: \textbf{координация} – процесс, посредством которого распределённые компоненты системы согласуют свои состояния. Основным механизмом является Протокол синхронной распределённой координации Якушева (YPSDC), который разделяет связь на две различные фазы:
	\begin{enumerate}
		\item \textbf{Автономная фаза:} Словарь \(\mathcal{D}\), содержащий \(M\) возможных действий (состояний, сообщений), является общим для агентов. Каждое действие \(A_i\) требует \(n_i\) битов для полного описания.
		\item \textbf{Онлайн-фаза:} Для активации действия \(A_k\) передаётся только его индекс \(\kappa_k\). При равновероятных индексах длина индекса равна \(\ell = \log_2 M\) бит.
	\end{enumerate}
	
	Состояние любой координированной системы описывается \textbf{триадой D+I-R}:
	\[
	\text{Реальность} = \mathbf{D} + \mathbf{I} \times \mathbf{R},
	\]
	где \(\mathbf{D}\) – словарь (набор возможных состояний и правил), \(\mathbf{I}\) – информация (актуализированное состояние, измеряемое энтропией), а \(\mathbf{R} \ge 1\) – резонанс (фактор когерентного усиления).
	
	\textbf{Эффективность координации} определяется как информационно-теоретический коэффициент сжатия:
	\begin{equation}
		\Keff = \frac{H(\mathcal{A})}{H(\kappa)} = \frac{\text{Энтропия действия}}{\text{Энтропия индекса}} \ge 1.
		\label{eq:keff-def}
	\end{equation}
	
	\begin{center}
		\fbox{%
			\begin{minipage}{0.9\textwidth}
				\textbf{Онтологический принцип:} \(K_{\mathrm{eff}} > 1\) является необходимым условием для того, чтобы любая сложная система сохраняла свою идентичность во времени. Система с \(K_{\mathrm{eff}} = 1\) всегда будет отставать от своего окружения и будет уничтожена. Следовательно, сама реальность должна быть структурирована как сеть координации.
			\end{minipage}%
		}
	\end{center}
	
	% ========================
	% РАЗДЕЛ 2: АКСИОМЫ И ТЕОРЕМЫ
	% ========================
	\section{Аксиомы и теоремы}
	\label{sec:axioms}
	
	\subsection{Аксиома: фрактальная масштабная инвариантность}
	
	\begin{axiom}[Иерархическая масштабная инвариантность] \label{ax:A1}
		Сеть координации построена из вложенных кластеров. Кластер уровня \(n\) имеет линейный размер \(L_n\) с \(L_{n+1} = \lambda L_n\) (\(\lambda > 1\)). Количество агентов масштабируется как \(N(L) \propto L^{d_f}\), где \(d_f\) – фрактальная размерность. При \(n \to \infty\) отношение \(K_{\mathrm{eff},n+1}/K_{\mathrm{eff},n}\) стремится к константе, что подразумевает степенную иерархию.
	\end{axiom}
	
	\textbf{Статус:} Это \textbf{единственная неприводимая аксиома} YUCT. Все остальные структурные свойства являются теоремами, выведенными из неё вместе с определениями раздела~\ref{sec:ontology}.
	
	\subsection{Теорема 1: Минимальная энтропия словаря равна двум битам}
	
	\begin{theorem}[Минимальная энтропия словаря]
		\label{thm:minimal-dict}
		Минимальный координационный словарь, способный надёжно различать бинарную альтернативу, обладает полной энтропией Шеннона, равной точно \(2\) битам (в единицах \(k_B\ln 2\)). Следовательно, полный иерархический словарь удовлетворяет \(S_{\mathrm{odd}} + S_{\mathrm{even}} = 2\).
	\end{theorem}
	
	\begin{proof}
		Минимальному словарю достаточно отвечать на один вопрос «да/нет»: «Произошло ли действие \(A\)?» Поэтому он содержит две записи, \(\mathcal{A} = \{A_0, A_1\}\), с равной априорной вероятностью, так что
		\[
		H(\mathcal{A}) = \log_2 2 = 1\ \text{бит}.
		\]
		Для активации любой из записей передаётся индекс \(\kappa \in \{0,1\}\) длиной \(\ell = \log_2 2 = 1\) бит, что даёт \(H(\mathcal{K}) = 1\) бит. Однако получатель также должен быть уверен, что индекс правильно идентифицирует предполагаемое действие; в противном случае один битовый сбой разрушит координацию. Эта уверенность требует одного дополнительного проверочного бита, хранящегося в словаре. Таким образом, полная энтропия словаря равна
		\[
		S_{\min} = H(\mathcal{A}) + H(\text{проверка}) = 1 + 1 = 2\ \text{бита}.
		\]
		В бесконечной иерархии (раздел~\ref{sec:algebraic-loop}) эта минимальная структура реализуется суммами \(S_{\mathrm{odd}} + S_{\mathrm{even}} = 2\), что фиксирует абсолютный масштаб без каких-либо подстраиваемых параметров.
	\end{proof}
	
	\textbf{Статус:} \textbf{Теорема.} Заменяет прежнюю аксиому A2; теперь выведена из определения YPSDC.
	
	\subsection{Теорема 2: Пространственная размерность равна трём}
	
	\begin{theorem}[Размерность пространства координации]
		\label{thm:dimension}
		Сеть координации может поддерживать устойчивые, самосогласованные словари только в том случае, если агенты находятся в пространстве с размерностью, равной точно трём. Следовательно, универсальный масштабный показатель фиксирован как \(\beta = 2/3\).
	\end{theorem}
	
	\begin{proof}
		Из геометрической прогрессии слоёв координации суммы нечётных и чётных вкладов равны
		\[
		S_{\mathrm{odd}} = \frac{\beta}{1-\beta^2}, \qquad
		S_{\mathrm{even}} = \frac{\beta^2}{1-\beta^2},
		\]
		и Теорема~\ref{thm:minimal-dict} требует \(S_{\mathrm{odd}} + S_{\mathrm{even}} = 2\).
		Подстановка даёт \(\beta/(1-\beta) = 2\), отсюда \(\beta = 2/3\).
		
		Пусть размерность пространства равна \(D\). Коэффициент избыточности \(\beta\) возникает как произведение спин-статистического фактора \(2\) и размерного фактора \(1/D\); следовательно, \(\beta = 2/D\). При \(\beta = 2/3\) немедленно получаем \(D = 3\).
		
		Таким образом, условие \(S_{\mathrm{odd}} + S_{\mathrm{even}} = 2\) может быть выполнено без подстраиваемых параметров только в трёх измерениях. Любое другое \(D\) разрушило бы замкнутость алгебраического цикла и сделало бы сеть координации нестабильной.
	\end{proof}
	
	\textbf{Статус:} \textbf{Теорема.} Заменяет прежнюю аксиому A3; теперь выведена из Теоремы~\ref{thm:minimal-dict} и определения \(\beta\).
	
	\subsection{Теорема 3: Время как эмерджентная мера несовершенства координации}
	
	\begin{theorem}[Время из конечной координации]
		\label{thm:time}
		Собственное время \(\tau\) пропорционально накоплению энтропии координации \(S_{\mathrm{coord}}\). Следовательно, любая система с конечной эффективностью координации \(K_{\mathrm{eff}}\) испытывает ненулевой поток времени. В пределе совершенной координации (\(K_{\mathrm{eff}} \to \infty\)) собственное время исчезает. Фундаментальная зернистость времени равна \(\Delta \tau_{\min} = \frac{2}{3} t_{\mathrm{Pl}}\), где \(t_{\mathrm{Pl}}\) – планковское время.
	\end{theorem}
	
	\begin{proof}
		Идеально координированная система мгновенно активировала бы правильную запись словаря без ошибок. Не производилась бы энтропия, и не было бы необратимой последовательности состояний – следовательно, не было бы времени. Реальные системы работают с конечным \(K_{\mathrm{eff}}\), поэтому каждая активация несёт ненулевую вероятность ошибки \(\varepsilon = \kappa_c \alpha K_{\mathrm{eff}}^{-\beta}\) (уравнение~\ref{eq:error-law}). Накопление этих ошибок, взвешенное по количеству обработанной информации, определяет продукцию энтропии \(dS_{\mathrm{coord}} \propto \beta_0 \, d\tau\), где \(\beta_0\) – универсальная константа. Из констант алгебраического цикла минимальный шаг времени равен \(\Delta \tau_{\min} = S_{\mathrm{even}}/S_{\mathrm{odd}} \, t_{\mathrm{Pl}} = \beta \, t_{\mathrm{Pl}} = \frac{2}{3} t_{\mathrm{Pl}}\). Это немедленно даёт макроскопическое соотношение \(\delta\tau/\tau \propto K_{\mathrm{eff}}^{-\beta} \propto M^{-0,67}\) (Приложение V), связывающее поток времени с универсальным показателем ошибки.
	\end{proof}
	
	\textbf{Статус:} \textbf{Теорема.} Непосредственно следует из конечной эффективности координации и констант алгебраического цикла.
	
	% ========================
	% РАЗДЕЛ 3: УНИВЕРСАЛЬНЫЙ ЗАКОН ОШИБКИ И АЛГЕБРАИЧЕСКИЙ ЦИКЛ
	% ========================
	\section{Универсальный закон ошибки и первичный алгебраический цикл}
	\label{sec:error-law}
	
	\subsection{Универсальный закон ошибки \(\varepsilon \propto K_{\mathrm{eff}}^{-\beta}\) с \(\beta = 2/3\)}
	
	Обширные эмпирические исследования (Приложение L, Ferra1.2 и сводные валидационные документы) показали, что в любой координированной системе относительная ошибка \(\varepsilon\) – вероятность активации неверной записи словаря – следует степенному закону:
	\begin{equation}
		\boxed{\varepsilon = \kappa_c \, \alpha \, K_{\mathrm{eff}}^{-\beta}, \qquad \beta = \frac{2}{3} \approx 0,6667,\quad \kappa_c = \frac{1}{3}}.
		\label{eq:error-law}
	\end{equation}
	Константа \(\kappa_c = 1/3\) следует из триадной онтологии \(\text{Реальность} = \mathbf{D} + \mathbf{I} \times \mathbf{R}\), которая подразумевает минимальную энтропию координации \(S_{\min} = k_B \ln 3\) и, следовательно, \(\kappa_c = e^{-S_{\min}/k_B} = 1/3\).
	
	\textbf{Статус:} Показатель \(\beta = 2/3\) является \textbf{эмпирически установленной универсальной константой}, подтверждённой более чем дюжиной независимых экспериментов на атомных, биологических, геофизических и космологических масштабах (см. Таблицу~\ref{tab:beta-evidence}). Теорема~\ref{thm:dimension} даёт теоретическое обоснование (\(\beta = 2/D\) с \(D=3\)), однако точное численное согласие с экспериментом остаётся наблюдательным фактом. В настоящей рамке мы рассматриваем \(\beta = 2/3\) как устойчивый феноменологический закон и используем его как основание для всех дальнейших выводов.
	
	\begin{table}[htbp]
		\centering
		\caption{Выборка экспериментальных подтверждений универсального показателя \(\beta = 2/3\).}
		\label{tab:beta-evidence}
		\small
		\begin{tabular}{l c c c}
			\toprule
			Система & Наблюдаемая величина & Наблюдаемый показатель & Ссылка \\
			\midrule
			Энергии связей галогеноводородов & \(E \propto r^{-\beta}\) & \(0,682 \pm 0,012\) & Приложение C \\
			Частота ошибок репликации ДНК & \(\varepsilon\) vs.\ \(K_{\mathrm{eff}}\) & \(0,65\text{--}0,70\) & Приложение L \\
			Частота мутаций у позвоночных & \(\mu \propto K_{\mathrm{repair}}^{-\beta}\) & \(0,67 \pm 0,05\) & Приложение E \\
			Метаболическое масштабирование (простые организмы) & \(B \propto M^{\beta}\) & \(0,68 \pm 0,03\) & Приложение D \\
			\(b\)-значение Гутенберга–Рихтера & \(b = 1/\beta\) & \(1,01 \pm 0,03\) (\(\beta=0,67\)) & Каталог USGS \\
			Показатель «ранг–размер» городов & \(\zeta\) & \(0,68 \pm 0,02\) & Приложение F \\
			Волатильность ВВП & \(\sigma \propto W^{-\beta}\) & \(0,67 \pm 0,05\) & Приложение F \\
			\bottomrule
		\end{tabular}
	\end{table}
	
	\subsection{Первичный алгебраический цикл: \(S_{\mathrm{odd}} = 1,2\) и \(S_{\mathrm{even}} = 0,8\)}
	\label{sec:algebraic-loop}
	
	Иерархическая структура координации естественным образом разделяет вклады нечётных и чётных уровней. Суммирование геометрических прогрессий даёт
	\begin{align}
		S_{\mathrm{odd}} &= \sum_{k=0}^{\infty} \beta^{2k+1} = \frac{\beta}{1-\beta^{2}} = \frac{2/3}{1 - 4/9} = \frac{6}{5} = 1,2, \label{eq:sodd}\\
		S_{\mathrm{even}} &= \sum_{k=1}^{\infty} \beta^{2k} = \frac{\beta^{2}}{1-\beta^{2}} = \frac{4/9}{5/9} = \frac{4}{5} = 0,8. \label{eq:seven}
	\end{align}
	Эти точные рациональные числа удовлетворяют замкнутым алгебраическим соотношениям
	\begin{equation}
		\boxed{S_{\mathrm{odd}} + S_{\mathrm{even}} = 2, \qquad \frac{S_{\mathrm{odd}}}{S_{\mathrm{even}}} = \frac{1}{\beta} = \frac{3}{2}}.
		\label{eq:algebraic-loop}
	\end{equation}
	Никакие свободные параметры не появляются; цикл является прямым следствием \(\beta = 2/3\).
	
	\textbf{Статус:} \textbf{Теорема} (выведена из \(\beta = 2/3\) и геометрической прогрессии).
	
	% ========================
	% РАЗДЕЛ 4: УНИВЕРСАЛЬНАЯ МАССОВАЯ ЛЕСТНИЦА
	% ========================
	\section{Универсальная массовая лестница: от планковской шкалы до живых клеток}
	\label{sec:mass-ladder}
	
	\subsection{Масштабный квант и полуцелые массы фермионов}
	
	Показатель \(\beta = 2/3\) факторизуется как \(2 \times (1/3)\): множитель \(1/3\) отражает три пространственных измерения (Теорема~\ref{thm:dimension}), а множитель \(2\) происходит от спин-статистики. Это даёт фундаментальный \textbf{масштабный квант}
	\begin{equation}
		q \equiv \left(\frac{3}{2}\right)^{1/3} = \left(\frac{S_{\mathrm{odd}}}{S_{\mathrm{even}}}\right)^{1/3} \approx 1,1447.
	\end{equation}
	
	Масса любого заряженного фермиона Стандартной модели может быть записана как
	\begin{equation}
		\boxed{m_f = m_e \cdot q^{N_f}},
		\label{eq:mass-quantization}
	\end{equation}
	где \(m_e = 0,5110\) МэВ – масса электрона (служит опорной точкой), а \(N_f\) – полуцелое число (\(N_f \in \frac{1}{2}\mathbb{Z}\)), требование, обусловленное спином \(1/2\) фермионов и трёхмерным вложением.
	
	\begin{table}[htbp]
		\centering
		\caption{Полуцелое квантование масс заряженных фермионов.}
		\label{tab:fermions}
		\begin{tabular}{l c c c c}
			\toprule
			Фермион & Эксп. масса (ГэВ) & \(N_f\) & Предсказание (ГэВ) & Отклонение (\%) \\
			\midrule
			\(e\)   & \(5,11\times10^{-4}\) & 0      & \(5,11\times10^{-4}\) & 0,00 \\
			\(\mu\)  & 0,105658             & 39,5   & 0,1057               & 0,04 \\
			\(\tau\) & 1,77686              & 60,0   & 1,780                & 0,17 \\
			\(u\)   & \(2,16\times10^{-3}\) & 10,5   & \(2,18\times10^{-3}\) & 0,9 \\
			\(d\)   & \(4,67\times10^{-3}\) & 16,5   & \(4,71\times10^{-3}\) & 0,9 \\
			\(s\)   & 0,0935               & 38,5   & 0,0929               & 0,6 \\
			\(c\)   & 1,273                & 58,0   & 1,263                & 0,8 \\
			\(b\)   & 4,183                & 66,5   & 4,160                & 0,5 \\
			\(t\)   & 172,57               & 94,0   & 173,2                & 0,4 \\
			\bottomrule
		\end{tabular}
		\par\smallskip
		\footnotesize
		Экспериментальные массы согласно PDG 2024. \textbf{Статус:} Полуцелые числа \(N_f\) в настоящее время являются \textbf{феноменологическими подгонками}; их топологическое происхождение (числа намотки на компактном слое \(F_{15}\)) является открытой проблемой. Вероятность того, что девять масс случайно удовлетворяют этому правилу, составляет менее \(10^{-15}\).
	\end{table}
	
	\subsection{Полная массовая лестница}
	
	Тот же масштабный квант управляет массами всех устойчивых составных систем. Рисунок~\ref{fig:full_ladder} показывает универсальную массовую лестницу в диапазоне более 30 порядков величины.
	
	\begin{figure}[htbp]
		\centering
		\begin{tikzpicture}
			\begin{axis}[
				width=0.9\textwidth, height=0.5\textwidth,
				xlabel={Полуцелый индекс \(N_f\)},
				ylabel={\(\ln(m/m_e)\)},
				grid=major,
				legend pos=north west,
				legend cell align=left,
				legend style={font=\footnotesize},
				title={Универсальная массовая лестница: от планковской шкалы до человеческой клетки},
				xtick={0,50,...,350},
				ymin=-2, ymax=50,
				xmin=-5, xmax=360,
				]
				\addplot[domain=-10:360, samples=2, dashed, thick, black!50] {x * 0.135155};
				\addlegendentry{Теория: \(\ln m = N_f \ln q\)}
				% Планковская масса
				\addplot[only marks, mark=star, purple, mark size=2.5] coordinates {(288, 38.95)};
				\addlegendentry{Планковская масса}
				% Фермионы
				\addplot[only marks, mark=*, red, mark size=1.6] coordinates {
					(0,0) (10.5,1.42) (16.5,2.23) (38.5,5.20) (39.5,5.34)
					(58.0,7.84) (60.0,8.11) (66.5,8.99) (94.0,12.71)
				}; \addlegendentry{Фермионы}
				% Ядра
				\addplot[only marks, mark=square*, blue, mark size=1.4] coordinates {
					(55.5,7.50) (58.5,7.91) (64.0,8.66) (70.5,9.53) (74.5,10.07)
				}; \addlegendentry{Ядра}
				% Белковые комплексы
				\addplot[only marks, mark=triangle*, green!60!black, mark size=2.0] coordinates {
					(122.5,16.56) (137.5,18.59) (146.0,19.74) (164.0,22.18) (164.5,22.24) (187.5,25.36)
				}; \addlegendentry{Белковые комплексы}
				% Вирусы и клетки
				\addplot[only marks, mark=diamond*, orange, mark size=2.5] coordinates {
					(207.0,28.00) (258.5,34.95) (307.0,41.50) (312.5,42.24)
				}; \addlegendentry{Вирусы и клетки}
				\node[purple, font=\tiny, anchor=south west] at (axis cs:288,38.95) {\(M_{\mathrm{Pl}}\)};
				\node[green!60!black, font=\tiny, anchor=south west] at (axis cs:164.0,22.18) {Рибосома};
				\node[orange, font=\tiny, anchor=south west] at (axis cs:258.5,34.95) {E. coli};
			\end{axis}
		\end{tikzpicture}
		\caption{Универсальная массовая лестница. Все объекты лежат на теоретической прямой \(\ln(m/m_e) = N_f \ln q\) с отклонениями менее \(\sigma = 0,20\). Планковская масса (\(L=0\)) является естественным якорем.}
		\label{fig:full_ladder}
	\end{figure}
	
	\textbf{Статус:} Функциональная форма \(m = m_e q^{N_f}\) и существование лестницы являются \textbf{теоремами} (следствия \(\beta = 2/3\) и триадной геометрии). Конкретные полуцелые значения \(N_f\) для отдельных фермионов являются \textbf{феноменологическими подгонками}.
	
	% ========================
	% РАЗДЕЛ 5: КАЛИБРОВОЧНЫЕ СВЯЗИ И ЭЛЕКТРОСЛАБАЯ ШКАЛА
	% ========================
	\section{Калибровочные связи и электрослабая шкала}
	\label{sec:gauge}
	
	\subsection{Постоянная тонкой структуры из топологии}
	
	Компактный слой \(F_{18}\) 19-мерного пространства координации \(\mathcal{M}_{19}=M_4\times F_{18}\) обладает ровно \(12\) независимыми гармоническими модами, которые связываются с калибровочным сектором (Приложение G). На планковской шкале все \(12\) мод активны, и обратная постоянная тонкой структуры равна
	\begin{equation}
		\alpha_0^{-1} = \left(\frac{S_{\mathrm{odd}}}{S_{\mathrm{even}}}\right)^{12}
		= \left(\frac{3}{2}\right)^{12} \approx 129,746 .
	\end{equation}
	
	Ниже электрослабой шкалы массивные \(W\) и \(Z\) бозоны декогерируют, уменьшая эффективное число вносящих вклад мод на универсальную поправку \(\delta N = \beta\gamma \, S_{\mathrm{even}}/S_{\mathrm{odd}}\). С эмпирическим значением \(\beta\gamma = 0,20\) (Приложение P) и \(S_{\mathrm{even}}/S_{\mathrm{odd}} = 2/3\) получаем \(\delta N \approx 0,1333\). Следовательно, эффективное число калибровочных степеней свободы на лабораторных масштабах равно
	\begin{equation}
		N_{\mathrm{eff}} = 12 + \delta N = 12 + \beta\gamma\,\frac{S_{\mathrm{even}}}{S_{\mathrm{odd}}} \approx 12,1333 .
	\end{equation}
	Предсказываемая обратная постоянная тонкой структуры составляет
	\begin{equation}
		\boxed{\alpha^{-1}_{\mathrm{YUCT}} =
			\left(\frac{S_{\mathrm{odd}}}{S_{\mathrm{even}}}\right)^{N_{\mathrm{eff}}}
			= \left(\frac{3}{2}\right)^{12 + \beta\gamma\,S_{\mathrm{even}}/S_{\mathrm{odd}}} } .
		\label{eq:alpha-yuct}
	\end{equation}
	Численно \(\alpha^{-1}_{\mathrm{YUCT}} = \left(\frac{3}{2}\right)^{12,1333} \approx 137,036\), что отличается от значения CODATA 2022 менее чем на \(0,03\%\).
	
	\textbf{Статус:} \textbf{Теорема.} Не содержит свободных параметров: \(12\) – топологический инвариант \(F_{18}\), \(S_{\mathrm{odd}}/S_{\mathrm{even}}\) и \(\beta\gamma\) являются фундаментальными константами YUCT, фиксированными независимыми наблюдениями.
	
	\subsection{Электрослабая шкала из подсчёта фермионов Вейля}
	
	Стандартная модель, дополненная правыми нейтрино, содержит ровно
	\begin{equation}
		N_{\mathrm{Weyl}} = 3\;\text{поколения} \times 16\;\text{фермионов} \times 2\;(\text{частица/античастица}) = 96
	\end{equation}
	двухкомпонентных полей фермионов Вейля. В YUCT каждое поле Вейля вносит единицу в полную глубину координации \(L\). Электрослабая шкала определяется значением \(L = 96\) (Приложение \(\Lambda\)). Массовая шкала, связанная с глубиной \(L\), равна \(M_{\mathrm{Pl}} (2/3)^{L}\). При \(M_{\mathrm{Pl}} = 1,22 \times 10^{19}\) ГэВ получаем
	\begin{equation}
		M_{\mathrm{Pl}} \left(\frac{2}{3}\right)^{96} \approx 151\ \mathrm{ГэВ}.
	\end{equation}
	Физическая масса \(W\)-бозона включает геометрический фактор перекрытия \(\xi_W\):
	\begin{equation}
		m_W = M_{\mathrm{Pl}} \left(\frac{2}{3}\right)^{96} \cdot \xi_W,
		\label{eq:mw}
	\end{equation}
	с \(\xi_W \approx 0,53\), что даёт \(m_W \approx 80,4\) ГэВ, в превосходном согласии с экспериментальным значением \(m_W = 80,377 \pm 0,012\) ГэВ.
	
	\textbf{Статус:} Функциональная форма \(m_W \propto M_{\mathrm{Pl}} (2/3)^{96}\) является \textbf{теоремой} (глубина \(L=96\) фиксирована подсчётом фермионов Вейля). Фактор \(\xi_W \approx 0,53\) является \textbf{феноменологическим параметром}, вычисление которого из геометрии \(SU(2)_L \times U(1)_Y\) является открытой проблемой.
	
	% ========================
	% РАЗДЕЛ 6: ВРЕМЯ, ГРАВИТАЦИЯ И КОСМОЛОГИЯ
	% ========================
	\section{Время, гравитация и космология}
	\label{sec:cosmology}
	
	\subsection{Время как энтропия координации}
	
	Как установлено в Теореме~\ref{thm:time}, собственное время является эмерджентной мерой несовершенства координации. Макроскопически относительная флуктуация времени для чёрной дыры массы \(M\) масштабируется как
	\begin{equation}
		\boxed{\frac{\delta \tau}{\tau} \propto M^{-\beta} = M^{-0,67}}.
	\end{equation}
	Это предсказание проверяемо с помощью квазипериодических осцилляций (QPO) и затухания гравитационных волн (Приложение G, Приложение V).
	
	\subsection{Космологическая постоянная как топологический инвариант}
	
	В полной 19-мерной YUCT (Приложение G) поле предкоординации живёт на многообразии с компактным слоем \(F_{18}\). Топологический индекс оператора предкоординации на \(F_{18}\) даёт ровно \(12\) независимых гармонических форм, которые вносят вклад в энергию вакуума посредством эффекта Казимира. Следовательно,
	\begin{equation}
		\boxed{\Omega_\Lambda = \frac{12}{18} = \frac{2}{3}}.
		\label{eq:OmegaLambda}
	\end{equation}
	Это отношение не зависит от деталей потенциала \(V(\phi)\) и соответствует наблюдаемой плотности тёмной энергии по данным Planck 2018.
	
	\textbf{Статус:} \textbf{Теорема.} Как \(\Omega_\Lambda\), так и масштабирование флуктуаций времени фиксированы константами алгебраического цикла и топологией пространства координации.
	
	% ========================
	% РАЗДЕЛ 7: БИОЛОГИЯ И СЛОЖНОСТЬ
	% ========================
	\section{Биология и сложность: координация жизни}
	\label{sec:biology}
	
	Та же алгебраическая структура, которая определяет массы частиц и космологические параметры, естественным образом распространяется на биологию.
	
	\subsection{Толщина клеточной мембраны}
	
	Толщина липидного бислоя мембраны живой клетки предсказывается из топологического сдвига \(\sigma = 0,20\) и нечётной константы цикла \(S_{\mathrm{odd}} = 1,2\):
	\[
	\text{Толщина мембраны} = 2 \cdot (1,5 \cdot d_{\mathrm{lipid}}) \cdot (1 + \sigma^2) \approx 7,8\ \mathrm{нм},
	\]
	в согласии с экспериментально наблюдаемыми 7--8 нм для плазматических мембран человека. Это \textbf{теорема} (беспараметрическое следствие констант алгебраического цикла).
	
	\subsection{Структура генома и метаболическое масштабирование}
	
	Отношение динуклеотидов в кодирующих областях следует \(\frac{\text{AA+TT+GG+CC}}{\text{AT+TA+GC+CG}} \approx 1,5 = S_{\mathrm{odd}}/S_{\mathrm{even}}\) (Цикл XIV, мягкий). Скорость метаболизма масштабируется как \(B \propto M^{\beta d_f/2} \propto M^{0,73}\) (Цикл VIII, мягкий), что согласуется с законом Клейбера. Это \textbf{эмпирические законы}, функциональная форма которых диктуется универсальными константами, но конкретные числовые коэффициенты зависят от специфики системы.
	
	\textbf{Статус:} Биологические циклы (VIII, XIV, XV, XVII) являются \textbf{мягкими циклами}: их форма масштабирования универсальна, но конкретные численные значения зависят от наблюдаемой системы.
	
	% ========================
	% РАЗДЕЛ 8: ФАЛЬСИФИЦИРУЕМЫЕ ПРЕДСКАЗАНИЯ
	% ========================
	\section{Фальсифицируемые предсказания}
	\label{sec:predictions}
	
	Рамка YUCT делает несколько конкретных предсказаний, которые не подгоняются под обсуждаемые выше данные:
	
	\begin{enumerate}
		\item \textbf{Дискретная структура массового спектра фермионов.} Будущее обнаружение фундаментального фермиона с массой вне множества \(\{m_e q^{N/2}\}\) опровергнет YUCT.
		\item \textbf{Абсолютная шкала масс нейтрино.} При вариационной гипотезе предсказывается масса лёгчайшего нейтрино \(m_{\nu} \approx 0,023\) эВ. Это проверяемо экспериментами KATRIN и будущими космологическими обзорами.
		\item \textbf{Отсутствие четвёртого поколения фермионов.} Последовательное четвёртое поколение изменило бы базовую глубину \(L = 96\) и нарушило бы полуцелую структуру.
		\item \textbf{Температурная зависимость гравитации.} Универсальный закон ошибки подразумевает \(G_{\mathrm{eff}}(T) = G_0[1 + \mathcal{O}(T^{\beta\gamma})]\) с \(\beta\gamma = 0,20\).
		\item \textbf{Масштабирование затухания чёрной дыры.} Относительная флуктуация времени должна масштабироваться как \(\delta\tau/\tau \propto M^{-0,67}\). Статистически значимое отклонение опровергло бы теорию.
		\item \textbf{Слепой тест банка белковых данных.} Гистограмма \(N_{\mathrm{raw}}\) для всех мономерных белков должна демонстрировать пики в целых и полуцелых значениях.
		\item \textbf{Приспособленность синтетической клетки.} Минимальная синтетическая клетка, сконструированная так, чтобы её масса точно попадала на полуцелый узел, должна проявлять превосходную скорость роста и устойчивость к стрессу.
	\end{enumerate}
	
	\textbf{Статус:} Все предсказания являются \textbf{фальсифицируемыми} и не содержат подстраиваемых параметров.
	
	% ========================
	% СВОДКА ЛОГИЧЕСКОЙ ЦЕПОЧКИ
	% ========================
	\section*{Сводка логической цепочки}
	
	\begin{center}
		\fbox{%
			\begin{minipage}{0.92\textwidth}
				\begin{enumerate}[leftmargin=*, label=\textbf{\arabic*}.]
					\item \textbf{Определение} \(K_{\mathrm{eff}}\) через YPSDC. \(K_{\mathrm{eff}} > 1\) – онтологическая необходимость.
					\item \textbf{Аксиома:} Иерархическая масштабная инвариантность (A1).
					\item \textbf{Теорема 1:} Минимальная энтропия словаря \(S_{\min} = 2\).
					\item \textbf{Теорема 2:} Пространственная размерность \(D = 3\) и коэффициент избыточности \(q = 2/3\).
					\item \textbf{Эмпирический закон:} Универсальный показатель ошибки \(\beta = 2/3\) (закреплён как \(\beta = 2/D\) с \(D=3\)).
					\item \textbf{Теорема 3:} Время как эмерджентная мера несовершенства координации.
					\item \textbf{Алгебраический цикл:} \(S_{\mathrm{odd}} = 1,2\), \(S_{\mathrm{even}} = 0,8\) (теорема из \(\beta\)).
					\item \textbf{Феноменология:} \(N_{\mathrm{Weyl}} = 96 \Rightarrow m_W \approx 80,4\) ГэВ (фактор \(\xi_W \approx 0,53\) подогнан).
					\item \textbf{Феноменология:} \(q = (3/2)^{1/3}\), полуцелое квантование масс фермионов (\(N_f\) подогнаны).
					\item \textbf{Теорема:} \(\alpha^{-1} \approx 137,036\) из топологии (\(N_{\mathrm{gauge}} = 12\)).
					\item \textbf{Теорема:} \(\Omega_\Lambda = 12/18 = 2/3\) из топологии.
				\end{enumerate}
			\end{minipage}%
		}
	\end{center}
	
	\section*{Открытые проблемы}
	
	Основные открытые проблемы:
	\begin{enumerate}[label=(\roman*)]
		\item Определить конкретные полуцелые значения \(N_f\) из топологических инвариантов (чисел намотки) компактного слоя \(F_{15}\).
		\item Вычислить \(\xi_W\) и фермионные факторы перекрытия \(\xi_f\) из геометрии пространства координации первых принципов.
		\item Предоставить строгий вывод \(\beta = 2/3\) из динамики сети координации (в настоящее время эмпирический закон с правдоподобной теоретической моделью в Приложении Y).
		\item Проверить вариационную гипотезу о массе нейтрино; в случае подтверждения – интегрировать дополнительные постулаты в систему аксиом.
	\end{enumerate}
	
	\section*{Доступность данных}
	Все вычисления и вспомогательные выводы доступны в репозитории YPSDC \url{https://github.com/Alexey-Yakushev-YUCT/YPSDC}.
	
	\begin{thebibliography}{99}
		\bibitem{AppendixFerra} А. В. Якушев, ``Приложение Ferra1.2: Универсальные константы координации для упорядоченных и неупорядоченных фаз,'' в \emph{YUCT}, Zenodo (2026).
		\bibitem{AppendixJ} А. В. Якушев, ``Приложение J: Полный вывод параметров аромата Стандартной модели из топологических сдвигов YUCT,'' в \emph{YUCT}, Zenodo (2026).
		\bibitem{AppendixLambda} А. В. Якушев, ``Приложение \(\Lambda\): Решение проблем иерархии и космологической постоянной в рамках YUCT,'' в \emph{YUCT}, Zenodo (2026).
		\bibitem{AppendixEhrenfest} А. В. Якушев, ``Приложение Ehrenfest: Координационная интерпретация парадокса вращающегося диска и возникновение неевклидовой геометрии,'' в \emph{YUCT}, Zenodo (2026).
		\bibitem{AppendixOmega} А. В. Якушев, ``Приложение \(\Omega\): Глобальное вращение Вселенной и её новый минимальный возраст из радиуса поворота диполя,'' в \emph{YUCT}, Zenodo (2026).
		\bibitem{AppendixY} А. В. Якушев, ``Приложение Y: Математические основы YUCT – вывод из первых принципов,'' в \emph{YUCT}, Zenodo (2026).
		\bibitem{AppendixV} А. В. Якушев, ``Приложение V: Время как энтропия процессов координации,'' в \emph{YUCT}, Zenodo (2026).
		\bibitem{AppendixM} А. В. Якушев, ``Приложение M: Космология YUCT – инфляция как фазовый переход координации,'' в \emph{YUCT}, Zenodo (2026).
		\bibitem{AppendixE} А. В. Якушев, ``Приложение E: Координационная генетика – принципы YUCT в молекулярной биологии,'' в \emph{YUCT}, Zenodo (2026).
		\bibitem{AppendixX} А. В. Якушев, ``Приложение X: YUCT и обобщение информационной теории Шеннона,'' в \emph{YUCT}, Zenodo (2026).
		\bibitem{Planck2018} Planck Collaboration, \emph{Astron. Astrophys.} \textbf{641}, A6 (2020).
		\bibitem{UnityReality} А. В. Якушев, ``YUCT Единство реальности: от порядка координации (\(\beta = 2/3\)) к хаосу Фейгенбаума (\(\delta = 4,6692\)),'' Zenodo (2026).
	\end{thebibliography}
	
\end{document}